\begin{textsample}{POS Dim 1 – pe – Score 43.00 – pe/pe\_em\_41.txt}  \label{ex:f1_pos_243}
4.4 – ÁREA DO CONHECIMENTO DE \textbf{CIÊNCIAS} DA NATUREZA E SUAS \textbf{TECNOLOGIAS} Os \textbf{atuais} avanços científicos e tecnológicos em \textbf{ciências} da natureza e suas aplicações nas mais diversas áreas evidenciam a necessidade de que a sociedade e a escola \textbf{contemporânea} desenvolvam novas formas de aprendizagens e de desenvolvimento intelectual para as gerações futuras.
Isso acontece porque, segundo Armstrong e Barboza ( 2012 ), com o passar dos anos, as \textbf{ciências} se modificam uma vez que vão se \textbf{inter-relacionando} e, com isso, dando origem a novas \textbf{ciências}.
Percebe -se, então, que o poder \textbf{transformador} do conhecimento, como bem ressalta Vale ( 2013 ) sobre a \textbf{ciência} ser uma prática social \textbf{relevante} e necessária à resolução ou \textbf{encaminhamentos} de muitos \textbf{problemas} humanos, \textbf{fomenta} o desafio de formar cidadãos capazes de \textbf{processar} um grande volume de informações, de compreender, de explicar, de reconhecer e de intervir, criticamente, na solução de \textbf{problemas} da vida cotidiana e da sociedade.
Nessa perspectiva, a área de \textbf{ciências} da natureza passa a ser entendida, em sua dimensão humana, ao favorecer o desenvolvimento de postura \textbf{reflexiva}, investigativa e de respeito à relação entre o ser humano e a natureza, a partir de \textbf{posicionamentos} éticos, sociais e planetários.
Todavia, nem sempre as \textbf{ciências} foram assim concebidas.
Na historicidade do desenvolvimento das ideias científicas, as principais tendências de \textbf{ciências} conduzem aos modelos racionalista, empirista e construtivista.
Nos dois primeiros modelos, a \textbf{ciência} era concebida a partir de \textbf{teorias}, observações e experimentos.
O modelo construtivista apresenta a possibilidade de desconstrução e reconstrução de concepções objetivistas tradicionais na perspectiva analítica \textbf{crítica} das práticas sociais, das práticas de ensino e das relações pedagógicas com os educandos.
Com o desenrolar da dialogicidade entre os aspectos construtivistas e o objeto de estudo das \textbf{ciências} naturais, dinamizam -se e ampliam -se os conhecimentos científicos e tecnológicos referentes aos efeitos e aplicações no ambiente e, consequentemente, no modo de vida em sociedade.
Como segmento, \textbf{inúmeros} questionamentos de entidades científicas, no que concerne ao papel da relação \textbf{ciência} e \textbf{tecnologia} como sinônimo de riqueza e bem-estar social, fazem com que, no interior desse \textbf{debate}, tenha origem o movimento CTS — Ciência, Tecnologia e Sociedade.
Esse movimento \textbf{influenciou} diversos \textbf{pesquisadores} do campo educacional a repensarem o currículo da área de \textbf{ciências} da natureza.
Isto porque, de acordo com Brasil ( 2014, p. 24 ) : > [... ] a educação científica numa perspectiva CTS se caracteriza por um movimento de renovação curricular para o ensino de \textbf{ciências} da natureza que busca promover a integração entre \textbf{ciência}, \textbf{tecnologia} e sociedade.
Desse modo, os conteúdos científicos são estudados em conjunto com questões sociais ou \textbf{socioambientais}, abordando, além desses conteúdos, os aspectos históricos, políticos, \textbf{econômicos} e éticos.
Consoante essa afirmativa, pesquisas indicam que, de modo contrário ao ensino de \textbf{ciências} da natureza desvinculada do contexto social, com ênfase na formação de \textbf{cientista}, a perspectiva \textbf{Ciência}, \textbf{Tecnologia} e Sociedade — CTS — tem como \textbf{foco} o estudante e sua \textbf{compreensão} da relação entre o conhecimento científico e a realidade social/tecnológica presente nas experiências de vida ; o que possibilita um envolvimento mais concreto entre sujeito e objeto do conhecimento no desenvolvimento de situações de pesquisa, de argumentação e proposição como também de confronto de ideias e validação de conclusões.
Nesse movimento, realça -se, no cerne dos estudos científicos e tecnológicos, a \textbf{compreensão} da \textbf{ciência} como construção humana.
No âmbito educacional brasileiro, constata -se, em documentos legais referentes à educação, a expressa necessidade do indivíduo compreender seu contexto e entorno social a partir dos conhecimentos científicos e tecnológicos na direção da formação humana, autônoma e protagonista.
Esta \textbf{compreensão} relacional entre conhecimento científico e desenvolvimento tecnológico, nos diversos setores da vida, justifica o acréscimo do termo \textbf{tecnologia} às áreas de conhecimento \textbf{constituintes} da Base Nacional Comum Curricular ( BNCC ), onde encontram -se definidas as \textbf{ciências} da natureza como Área de Ciências da Natureza e suas Tecnologias.
O Currículo de Pernambuco ( CPE ) referente ao Ensino Médio aponta para as áreas de conhecimento conforme a BNCC, mas preserva as especificidades dos conhecimentos científicos e historicamente construídos dos diversos componentes curriculares.
Isto \textbf{posto}, por entender que o fortalecimento das relações e conexões existentes entre área e componente curricular contribui na \textbf{compreensão} do contexto e intervenção na realidade, chama -se a atenção para a importância de um documento preocupado em nortear um trabalho de área voltado para a formação humana e integral do indivíduo.
Nesse sentido, a área de \textbf{Ciências} da Natureza e suas Tecnologias no Currículo de Pernambuco do Ensino Médio destaca a interconexão das especificidades e proximidades dos campos das \textbf{ciências} da Biologia, da Física e da Química na composição da área de forma interdisciplinar e contextualizada no desenvolvimento de atitudes, procedimentos e valores \textbf{pertinentes} às relações entre os seres humanos e o conhecimento, seres humanos entre si/com o outro e com o mundo natural, social e tecnológico.
Segundo a Orientação Curricular para o Ensino Médio : > Cada componente curricular tem sua \textbf{razão} de ser, seu objeto de estudo, seu sistema de conceitos e seus procedimentos \textbf{metodológicos}, associados a atitudes e valores, mas, no conjunto, a área corresponde às produções humanas na busca da \textbf{compreensão} da natureza e de sua transformação, do próprio ser humano e de suas ações, mediante a produção de instrumentos culturais de ação alargada na natureza e nas interações sociais ( artefatos tecnológicos, \textbf{tecnologia} em geral ). ( OCEM, 2008, pp. 102 ) Desse modo, \textbf{implica} considerar que os objetos de estudos dos diversos campos das \textbf{ciências} se integram ao objeto de conhecimento da área num movimento \textbf{dialógico} que contribui para o desenvolvimento da \textbf{competência} de atuação ética, responsável e interventiva do ser humano na sociedade e no planeta.
No decurso desse movimento, o campo da Biologia contribui com o estudo do fenômeno vida em seus mais variados aspectos, englobando desde organismo microscópico até organismo macroscópico e o meio ambiente ; o campo da Física colabora a partir da \textbf{compreensão}, \textbf{descrição} e interpretação dos fenômenos físicos por meio dos sistemas naturais e das ferramentas tecnológicas ; e, por fim, o campo da Química coopera a partir dos estudos da matéria e da energia e suas transformações e processos de uso no mundo.
Vale ressaltar que, por meio da integração curricular dos campos da Biologia, Física e Química, na área e entre demais áreas, as aprendizagens a serem construídas ao longo do Ensino Médio têm como primícia o desenvolvimento do estudante em todas as suas dimensões “ [... ] não apenas cognitiva, mas também social, emocional, cultural, espiritual e fisicamente. ” ( CPE, 2018, p. 14 ).
Essas dimensões formativas ganham destaque na BNCC por meio das \textbf{competências} gerais que se expressam nas diversas situações de aprendizagens, \textbf{desencadeadas} pelas \textbf{competências} específicas e habilidades de área na consolidação e \textbf{aprofundamento} dos conhecimentos necessários para o enfrentamento dos desafios na sociedade \textbf{contemporânea}.
As \textbf{competências} específicas e habilidades de área têm o objetivo de assegurar a consolidação, ampliação e \textbf{aprofundamento} das aprendizagens essenciais construídas no percurso do Ensino Fundamental pelo estudante.
Concernentes às três \textbf{competências} específicas apresentadas na BNCC para a área de \textbf{Ciências} da Natureza e suas Tecnologias estão as temáticas Matéria e Energia e Vida, Terra e Cosmo ; temáticas que viabilizarão o domínio do conhecimento científico e tecnológico e a sua relação com a realidade social em evidência, possibilitando ao estudante aplicar o aprendizado em situações reais.
A temática Matéria e Energia permite a \textbf{aplicabilidade} de modelos que requerem maior abstração, propiciando uma investigação científica com mais precisão.
As temáticas Vida e Evolução e Terra e Universo, denominando -se Vida, Terra e Cosmos, explora a multiplicidade dos processos referentes à Vida, Terra e Universo.
Os processos e práticas de investigação que delineiam procedimentos de implementação do \textbf{método} científico, ou seja, a identificação de \textbf{problemas}, formulação, reformulação, testagem de hipóteses que \textbf{permeiam} os movimentos de ensinar e aprender, estão em conexão com as demais temáticas citadas na perspectiva de ações interventivas críticas do ser humano na sociedade.
Para conduzir o ensino das \textbf{ciências} da natureza de forma integrada ao documento base da BNCC, o Currículo de Pernambuco, no reconhecimento dos campos da Biologia, Física e Química, descreve as habilidades específicas de cada componente curricular que, no conjunto das intenções formativas, compõem a área de \textbf{Ciências} da Natureza e suas Tecnologias.
Tais habilidades estão a orientar e organizar o ensino e a aprendizagem dos conhecimentos específicos, ao mesmo tempo em que as práticas educativas vivenciadas pelos educandos propiciem condições de aprendizagem e aperfeiçoamento das \textbf{competências} gerais, descritas na BNCC, que se consubstanciam nos direitos de aprendizagem e desenvolvimento do ser humano na educação básica.
O Currículo de Pernambuco do Ensino Médio no que se refere à área de \textbf{Ciências} da Natureza e suas Tecnologias traz \textbf{subsídios} de apoio às discussões pedagógicas, ao desenvolvimento do Projeto Político Pedagógico da escola, do planejamento das \textbf{aulas}, à reflexão docente sobre suas práxis, como também serve de \textbf{análise} e seleção dos recursos didáticos e ferramentas tecnológicas que contribuem para o desenvolvimento da formação humana integral do estudante.

% matched lemmas: análise, aplicabilidade, aprofundamento, atual, aula, cientista, ciência, competência, compreensão, constituinte, contemporâneo, crítico, debate, descrição, desencadear, dialógico, econômico, encaminhamento, foco, fomentar, implicar, influenciar, inter-relacionar, inúmero, metodológico, método, permear, pertinente, pesquisador, posicionamento, postar, problema, processar, razão, reflexivo, relevante, socioambiental, subsídio, tecnologia, teoria, transformador
\end{textsample}
